\abstract

\textbf{Ключевые слова:} ОБНАРУЖЕНИЕ АНОМАЛИЙ, KUBERNETES, АУДИТ ЛОГИ, НЕЙРОННЫЕ СЕТИ, МАШИННОЕ ОБУЧЕНИЕ, КЛАССИФИКАТОР, ИНФОРМАЦИОННАЯ БЕЗОПАСНОСТЬ, КОНТЕЙНЕРНЫЕ СРЕДЫ.

Целью выпускной квалификационной работы является разработка средства обнаружения аномалий в Kubernetes-кластере на основе методов нейронных сетей, обеспечивающего высокую точность выявления аномальных событий.

В процессе выполнения работы были рассмотрены существующие подходы к обнаружению аномалий в информационных системах, а также методы анализа журналов аудита (audit-логов) в Kubernetes-кластерах. Проведён обзор современных методов машинного обучения и глубокого обучения, применяемых для задач выявления отклоняющегося поведения в последовательностях событий.

В ходе работы был выполнен анализ исходных данных, представляющих собой журналы аудита Kubernetes-кластера. На основе полученных данных были разработаны, обучены и исследованы модели обнаружения аномалий с использованием нейросетевых архитектур, включая автоэнкодеры и их рекуррентные модификации. Особое внимание уделено модели на основе LSTM-автоэнкодера (LSTM-AE), способной учитывать временную структуру последовательностей.

Оценка качества моделей проводилась с использованием метрик классификации, при этом в качестве основной метрики была выбрана площадь под PR-кривой (AUC_PR), как наиболее информативная в условиях несбалансированных данных. В результате экспериментов установлено, что модель LSTM-AE демонстрирует наилучшие показатели качества: значение AUC_PR составило 0,99962, precision — 0,93097, recall — 1,00000 и F1-мера — 0,96425, что подтверждает её высокую эффективность для решения поставленной задачи.

Далее было разработано программное обеспечение, реализующее процесс предварительной обработки audit-логов, формирования признаков, обучения и применения модели для выявления аномалий в Kubernetes-кластере. Предложенное средство обеспечивает автоматизированный анализ событий и позволяет своевременно обнаруживать отклонения от нормального поведения в контейнерной среде.

В рамках работы также рассмотрены вопросы практической применимости разработанного средства, включая оценку его экономической целесообразности и соответствие требованиям действующих нормативных правовых актов Российской Федерации.